
%Early childhood, spanning the first five years of life, is a critical period for brain maturation, yet the developmental trajectory of structure--function coupling (SFC) and its relationship to the sensorimotor--association (SA) axis during this window remain poorly understood. SFC reflects experience-dependent plasticity in healthy development and its disruption in neuropsychiatric conditions. Here, we document pronounced remodeling of SFC in early childhood that aligns with cortical patterns of excitation--inhibition balance (indexed by the Hurst exponent) and intracortical myelination (indexed by the T1w/T2w ratio). 
%Across the cortex, SFC increased by $\sim$23\%, myelination by $\sim$15\%, and the mean Hurst exponent decreased by $\sim$6\%, with regional SFC maturation covarying with both myelination and Hurst exponent.We analyzed multi-session multimodal MRI from more than \rev{200} children from birth to age 5, acquired in more than 1,000 imaging sessions.Structural and functional connectivity, as well as SFC, all strengthened along the SA axis, with association  regions showing steeper gains than sensorimotor regions. 
Early childhood, spanning the first five years of life, is a critical window of cortical maturation, yet the developmental trajectory of structure-function coupling (SFC) and its biological drivers remains poorly understood. Here, we used multimodal MRI to map structural connectivity (SC), functional connectivity (FC), SFC, intracortical myelination (T1w/T2w), and excitation-inhibition-related dynamics (Hurst exponent) at vertex resolution on the cortical surface from birth to 5 years, with age modeled at fine day-scale resolution. We then integrated these analyses across complementary spatial scales, including vertex-wise maps, Schaefer atlas systems, and the continuous sensorimotor-association (SA) axis. Across the cortex, SFC increased during early childhood, and this increase indicates progressively stronger integration across the brain from 0 to 5 years of age. In parallel, increasing myelination and shifts in E/I-related dynamics tracked the emergence of SFC, indicating that local microstructure and neural dynamics jointly constrain how structural and functional co-development is expressed as coupling. Within this broader cortical developmental pattern, the SA axis organized regional differences in SC, FC, and SFC maturation. Together, these findings provide a normative reference for probing atypical neurodevelopment and neuropsychiatric risk.

